%\section{Modelos de Lenguaje}
\phantomsection\addcontentsline{toc}{section}{Modelos de Lenguaje}
\section*{Modelos de Lenguaje}

Un Modelo de Lenguaje, o LM por sus siglas en inglés, es un sistema probabilístico que asigna
una probabilidad a una secuencia de palabras, o bien, predice la probabilidad de aparición
de la siguiente palabra dado un contexto previo. Formalmente, un modelo de lenguaje define
una distribución sobre secuencias $\omega_1, \omega_2, \ldots, \omega_n$, de la forma:

\[
    P(w_1, w_2, \ldots, w_n) = \prod_{i=1}^n P(w_i \mid w_1, w_2, \ldots, w_{i-1}),
\]

Donde cada término $P(\omega_i \mid \omega_{<i})$ representa la probabilidad condicional de
la palabra $\omega_i$ dado su historia previa $\omega_{<i}$.

Los modleos de lenguaje tienen una larga tradición en el procesmaiento de lenguaje natural. 
Los primeros enfoques se basaron en modelos de n-gramas, que utilizan la suposición de Markov, i.e.
la probabilidad de una palabra depndiente únicamente de las últimas $n-1$ palabras de contexto. Por
ejemplo, un modelo de bigramas:

\[
    P(\omega_i | \omega_{<i}) \approx P(\omega_i | \omega_{<i})
\]

Estos modelos, aunque eran sencillos y computacionalmente eficientes, presentaban limitaciones claras
al no poder capturar depndencias de largo alcance, además de sufrir de "esparsidad", es decir,
muchas secuencias posibles nunca aparecían en los datos de entrenamiento. Es así como el 
desarrollo de técnicas de suavizado como Laplace o Kneser-Ney permitieron mejorar el desempeño de
los LM, pero no se resolvió por completo en cuanto a escasez de datos.

Tras el auge del Deep Learning demostrado en la última década, los modelos de redes neuronales 
reemplazaron a los modelos de n-gramas, pues estos eran capaces de aprender representaciones
continuas y modelar depndencias mucho más largas de texto. Al poco tiempo, los Transformers
cambiaron el paradigma y permitieron ampliar dicha capacidad, aprovechando del mecanismo de
self-attention para capturar relaciones globales en la secuencia, y a entrenar a muy graned escala
llegnado a lo que ahora conocemos como Large Language Models o LLM, como lo son GPT, BERT, LLama y 
todo el abanico de sucesores. 

El propósito central de un LM es estimar la probabilidad de una secuencia de palabras, locual lo convierte 
en una herramienta súmamente útil para una enorme cantida de tareas. Sin embargo, dado que todo LM 
asigna distribuciones de probabilidad sobre textom su calidad se mide en qué tan bien predice 
secuencias reales del idioma. El criterio intrínseco más utilizado para esta tarea de evaluación es
la perplejidad. 

%\section{Perplejidad}
\phantomsection\addcontentsline{toc}{section}{Perplejidad}
\section*{Perplejidad}

La Perplejidad (PPL) es la métrica estándar para evaluar modelos de lenguaje, especialmente
en tareas de predicción de secuencias. De manera intuitiva, PPL mide qué tan sorprendido
está un modelo frente a un conjunto de validación, cuanto mejor predice las palabras
reales que aparecen en el texto, menor será su perplejidad. 

Formalmente, para un conjunto de prueba $W = \omega_1, \omega_2,..., \omega_N$, la perplejidad
se define como:

\[
    \text{PPL}(W) = \exp \Big(- \frac{1}{N} \sum_{i=1}^{N} \text{log} p(\omega_i | \omega_{<i})\Big)
\]

La expresión anterior equivale a:

\[
    \text{PPL}(W) = \exp \Big(H(p,q) \Big)
\]

Donde $p(\omega_i | \omega_{<i})$ es la probabilidad asignada por el modelo para la palabra
$\omega_i$, dado el contexto anterior. 

Esta definición puede interpretarse como la media geométrica inversa de las probabilidades 
de que el modleo asigna a las palabras correctas. En consecuencia:

\begin{itemize}
    \item Una PPL baja implica que el modelo asigna alta probabilidad a las palabras reales, por
    lo que hay una mejor capacidad de predicción.
    \item Una PPL alta implica que el modelo se sorprende con frecuencia, entonces predice
    mal el texto.
\end{itemize}

Así entonces, la perplejidad equivale al número promedio de opciones equiprobables entre las
que el modelo ``cree'' tener que elegir en cada paso. Por ejemplo, una PPL = 3 significa que
el modelo se comporta como si en promedio tuviera que escoger entre tres palabras posibles
en cada posición.

La perplejidad está directamente relacionada con la entropía cruzada de la distribución
real frente a la estimada por el modleo. Es decir, minimizar la perplejidad equivale a 
maximizar la probabilidad del conjunto de validación. Por esta razón, se utiliza ampliamente
para comparar modelos de lenguaje entrenados con distintos algoritmos o configuraciones. 

Es importante destacar que:

\begin{itemize}
    \item No siempre correlaciona perfectamente con la calidad extrínseca de tareas como traducción automática o generación de texto, donde aspectos de coherencia semántica y adecuación estilística son también relevantes.
    \item La comparación entre modelos requiere que se use el mismo vocabulario y corpus de prueba; de lo contrario, las cifras no son conmensurables.
    \item Modelos con perplejidad baja pueden generar texto gramatical pero poco creativo (alta repetición), mientras que temperaturas de muestreo altas pueden producir lo opuesto: diversidad pero con mayor perplejidad.
\end{itemize}

\subsection{PPL en modelos de Redes Neuronales}

En arquitecturas modernas como RNN, LSTM, GRU o Transformers, el cálculo de la perplejidad se
implementa de manera directa a partir de la función de \textit{cross-entropy}. dado
que durante el entrenamiento se minimiza la entropía cruzada entre las etiquetas reales y la
distribución del modelo, evaluar la perplejidad en el conjunto de validación es equivalente a
medir el valor exponencial de dicha pérdida. 

